

\section{Results} \label{sec:results} % 17-Jul-2026 retraining
For each schema, four RL algorithms are evaluated with no attention (baseline) and four attention mechanisms -- Nadaraya--Watson, temporal, self-attention, and multi-head.

\subsection{Main performance tables and focused statistical tests}
\label{sec:hypothesis_tests}

Tables~\ref{tab:sit_main_results} and \ref{tab:ieee_main_results} report mean $\pm$ standard deviation of evaluation score. Table~\ref{tab:rf_vs_ppo_tests} gives the focused one-tailed Welch tests for the claim \textsc{REINFORCE} $>$ \textsc{PPO}.

\begin{table}[htbp]
	\centering
	\caption{SIT schema: mean $\pm$ std. dev. evaluation score.}
	\label{tab:sit_main_results}
	\begin{tabular}{l l l l l l}
	\toprule
	\textbf{RL Algo} & \textbf{Baseline} & \textbf{Nadaraya--Watson} & \textbf{Temporal} & \textbf{Self-attention} & \textbf{Multi-head}\\
	\midrule
	A2C & 0.2456 $\pm$ 0.0043 & 0.2456 $\pm$ 0.0043 & 0.2456 $\pm$ 0.0043 & 0.3130 $\pm$ 0.2040 & 0.2456 $\pm$ 0.0043\\
	DQN & 0.2456 $\pm$ 0.0043 & 0.2456 $\pm$ 0.0043 & 0.2456 $\pm$ 0.0043 & 0.2456 $\pm$ 0.0043 & 0.2456 $\pm$ 0.0043\\
	PPO & 0.6313 $\pm$ 0.3347 & 0.7132 $\pm$ 0.3172 & 0.5771 $\pm$ 0.3337 & 0.6089 $\pm$ 0.3305 & 0.8004 $\pm$ 0.2694\\
	REINFORCE & 0.4907 $\pm$ 0.3108 & 0.8646 $\pm$ 0.2032 & 0.6406 $\pm$ 0.3355 & 0.6997 $\pm$ 0.3171 & 0.7917 $\pm$ 0.2738\\
	\bottomrule
	\end{tabular}
\end{table}

\begin{table}[htbp]
	\centering
	\caption{IEEE schema: mean $\pm$ std. dev. evaluation score.}
	\label{tab:ieee_main_results}
	\begin{tabular}{l l l l l l}
	\toprule
	\textbf{RL Algo} & \textbf{Baseline} & \textbf{Nadaraya--Watson} & \textbf{Temporal} & \textbf{Self-attention} & \textbf{Multi-head}\\
	\midrule
	A2C & 0.6880 $\pm$ 0.3098 & 0.4322 $\pm$ 0.2753 & 0.3565 $\pm$ 0.2496 & 0.2729 $\pm$ 0.1281 & 0.7424 $\pm$ 0.2695\\
	DQN & 0.6696 $\pm$ 0.3002 & 0.2360 $\pm$ 0.0016 & 0.2360 $\pm$ 0.0016 & 0.2360 $\pm$ 0.0016 & 0.2878 $\pm$ 0.1811\\
	PPO & 0.8434 $\pm$ 0.1604 & 0.8343 $\pm$ 0.1867 & 0.7608 $\pm$ 0.2641 & 0.8044 $\pm$ 0.2174 & 0.8232 $\pm$ 0.1870\\
	REINFORCE & 0.7496 $\pm$ 0.2606 & 0.8008 $\pm$ 0.2186 & 0.8848 $\pm$ 0.0737 & 0.8860 $\pm$ 0.0708 & 0.8559 $\pm$ 0.1443\\
	\bottomrule
	\end{tabular}
\end{table}

\begin{table}[htbp]
	\centering
	\caption{Focused one-tailed Welch tests for \textsc{REINFORCE} $>$ \textsc{PPO} at $\alpha=0.05$.}
	\label{tab:rf_vs_ppo_tests}
	\begin{tabular}{l l r r l}
	\toprule
	\textbf{Schema} & \textbf{Setting} & \textbf{t-stat} & \textbf{p-value} & \textbf{Supported?}\\
	\midrule
	SIT & Baseline (no attention) & -2.9224 & 0.9980 & No\\
	SIT & With attention (pooled) & 3.1941 & 0.0007 & Yes\\
	IEEE & Baseline (no attention) & -1.3711 & 0.9100 & No\\
	IEEE & With attention (pooled) & 1.9865 & 0.0243 & Yes\\
	\bottomrule
	\end{tabular}
\end{table}

\paragraph{Result interpretation:}
Without attention, \textsc{PPO} has the highest mean in both schemas (SIT: 0.6313; IEEE: 0.8434). With attention, the best overall mean shifts to \textsc{REINFORCE} (SIT: 0.7492; IEEE: 0.8569). This is consistent with the focused tests in Table~\ref{tab:rf_vs_ppo_tests}: the claim \textsc{REINFORCE} $>$ \textsc{PPO} is unsupported at baseline and supported when attention is enabled.

\paragraph{Mechanism-level synthesis for \textsc{REINFORCE} and \textsc{PPO}:}
Table~\ref{tab:best_mech_ppo_rf} summarizes the best attention mechanism for \textsc{REINFORCE} and \textsc{PPO} in each schema using mean $\pm$ standard deviation. In SIT, \textsc{REINFORCE} is strongest with Nadaraya--Watson (0.8646 $\pm$ 0.2032), while \textsc{PPO} is strongest with multi-head attention (0.8004 $\pm$ 0.2694). In IEEE, \textsc{REINFORCE} is strongest with self-attention (0.8860 $\pm$ 0.0708), while \textsc{PPO} is strongest with Nadaraya--Watson (0.8343 $\pm$ 0.1867).

\begin{table}[htbp]
\centering
	\caption{Best attention mechanism by schema for \textsc{REINFORCE} and \textsc{PPO} (mean $\pm$ std. dev. evaluation score).}
	\label{tab:best_mech_ppo_rf}
	\begin{tabular}{l l l l}
	\toprule
	\textbf{Schema} & \textbf{Algorithm} & \textbf{Best attention mechanism} & \textbf{Eval. score (mean $\pm$ std. dev.)}\\
	\midrule
	SIT & \textsc{REINFORCE} & Nadaraya--Watson & 0.8646 $\pm$ 0.2032\\
	SIT & \textsc{PPO} & Multi-head & 0.8004 $\pm$ 0.2694\\
	IEEE & \textsc{REINFORCE} & Self-attention & 0.8860 $\pm$ 0.0708\\
	IEEE & \textsc{PPO} & Nadaraya--Watson & 0.8343 $\pm$ 0.1867\\
	\bottomrule
	\end{tabular}
\end{table}

% $$$ Model performance plots
% PPO_MH_IEEE_01_04
% PPO_MH_SIT_8_12
% REINFORCE_MH_IEEE_04_01
% REINFORCE_MH_SIT_12_10

\subsection{Ablation study: attention versus no-attention} \label{sec:ablation_results} % 17-Jul-2026 retraining

Table~\ref{tab:ablation_condensed} reports ablation study analysis. It includes a focussed algorithm-specific view for \textsc{REINFORCE} alongside the pooled comparisons.

\begin{table}[htbp]
	\centering
	\caption{Condensed ablation study analysis (one-tailed Wel t-tests).}
	\label{tab:ablation_condensed}
	\begin{tabular}{l l r r r}
		\toprule
		\textbf{Schema} & \textbf{Comparison} & \textbf{No attn.} & \textbf{Attention} & \textbf{$p$ (1-tail)}\\
		\midrule
		SIT & \textsc{REINFORCE} only & 0.4907 & 0.7492 & 0.0000\\
		IEEE & \textsc{REINFORCE} only & 0.7496 & 0.8569 & 0.0440\\
		SIT & PPO + \textsc{REINFORCE}, pooled & 0.5610 & 0.7120 & 0.0000\\
		IEEE & PPO + \textsc{REINFORCE}, pooled & 0.7965 & 0.8313 & 0.1757\\
		SIT & All algorithms, pooled & 0.4033 & 0.4830 & 0.0000\\
		IEEE & All algorithms, pooled & 0.7377 & 0.5906 & 1.0000\\
		\bottomrule
	\end{tabular}
\end{table}

For \textsc{REINFORCE}, attention significantly improves over the no-attention baseline in both schemas (SIT: 0.4907 to 0.7492, $p=0.0000$; IEEE: 0.7496 to 0.8569, $p=0.0440$), supporting attention interoperability across schemas for this algorithm. The PPO+\textsc{REINFORCE} subset follows the same direction: SIT improves significantly (0.5610 to 0.7120, $p=0.0000$), and IEEE trends upward, though the gain is not significant at $\alpha=0.05$ (0.7965 to 0.8313, $p=0.1757$). Pooled across all four algorithms, SIT again improves significantly (0.4033 to 0.4830, $p=0.0000$); the IEEE all-algorithm pooled comparison shows a different pattern, reflecting the already-strong no-attention baselines of \textsc{A2C} and \textsc{DQN} on IEEE (Table~\ref{tab:ieee_main_results}). Overall, attention interoperability across schemas is clearest for \textsc{REINFORCE}, holds more weakly for the PPO+\textsc{REINFORCE} subset, and remains schema-dependent once pooled across all four algorithms.

Figures~\ref{fig:analysis_reports}--\ref{fig:ablation_ppo_rf} provide visualizations of the summary analysis, corresponding to Tables~\ref{tab:sit_main_results}--\ref{tab:ablation_condensed}.

\begin{figure}[htbp!]
	\centering
	\begin{subfigure}{0.9\textwidth}
		\centering
		\includegraphics[width=\linewidth]{Analysis_Report_IEEE.png}
		\caption{IEEE analysis report.}
	\end{subfigure}
	
	\begin{subfigure}{0.9\textwidth}
		\centering
		\includegraphics[width=\linewidth]{Analysis_Report_SIT.png}
		\caption{SIT analysis report.}
	\end{subfigure}
	\caption{Overall analysis reports.}
	\label{fig:analysis_reports}
\end{figure}

\begin{figure}[t]
	\centering
	\begin{subfigure}[t]{\textwidth}
		\centering
		\includegraphics[width=\linewidth]{Ablation_Study_Algo_IEEE.jpg}
		\caption{IEEE ablation by algorithm.}
	\end{subfigure}

	\begin{subfigure}[t]{\textwidth}
		\centering
		\includegraphics[width=\linewidth]{Ablation_Study_Algo_SIT.jpg}
		\caption{SIT ablation by algorithm.}
	\end{subfigure}
	\caption{Algorithm-level ablation summaries.}
	\label{fig:ablation_algo}
\end{figure}

\begin{figure}[t]
	\centering
	\begin{subfigure}[t]{0.9\textwidth}
		\centering
		\includegraphics[width=\linewidth]{Ablation_Study_IEEE_All_Algos.jpg}
		\caption{IEEE: all algorithms pooled.}
	\end{subfigure}
	
	\begin{subfigure}[t]{0.9\textwidth}
		\centering
		\includegraphics[width=\linewidth]{Ablation_Study_SIT_All_Algos.jpg}
		\caption{SIT: all algorithms pooled.}
	\end{subfigure}
	\caption{All-algorithm pooled ablation comparisons.}
	\label{fig:ablation_all}
\end{figure}

\begin{figure}[t]
	\centering
	\begin{subfigure}[t]{0.9\textwidth}
		\centering
		\includegraphics[width=\linewidth]{Ablation_Study_IEEE_PPO_REINFORCE.jpg}
		\caption{IEEE: PPO + \textsc{REINFORCE}.}
	\end{subfigure}
	
	\begin{subfigure}[t]{0.9\textwidth}
		\centering
		\includegraphics[width=\linewidth]{Ablation_Study_SIT_PPO_REINFORCE.jpg}
		\caption{SIT: PPO + \textsc{REINFORCE}.}
	\end{subfigure}
	\caption{Ablation comparisons focused on PPO and \textsc{REINFORCE}.}
	\label{fig:ablation_ppo_rf}
\end{figure}

\subsection{Schema portability and attention sensitivity}
The same MDP and reward logic run across both schemas, supporting environment portability at the design level. However, performance patterns are different: SIT shows strong gains for selected attention setups (especially \textsc{REINFORCE} with Nadaraya--Watson and \textsc{PPO} with multi-head), while IEEE shows a narrower separation among top policy-gradient variants and weaker pooled ablation gains. This supports schema-aware model selection rather than direct policy transfer.

\subsection{Model sizes}
As a practical aside, the trained \textsc{REINFORCE} policies are also small in serialized form in these experiments (Table~\ref{tab:model_size_summary}). This table is not a hardware-deployment study; it only reports saved model sizes.

\begin{table}[htbp]
\centering
\caption{Average serialized policy size (kB), averaged over models within each schema.}
\label{tab:model_size_summary}
\begin{tabular}{l r r}
\toprule
	\textbf{Algorithm} & \textbf{SIT avg. (kB)} & \textbf{IEEE avg. (kB)}\\
	\midrule
	A2C & 293.55 & 292.54\\
	DQN & 419.73 & 417.72\\
	PPO & 169.90 & 169.38\\
	\textbf{REINFORCE} & \textbf{169.47} & \textbf{168.95}\\
\bottomrule
\end{tabular}
\end{table}

\section{Discussion}
\label{sec:discussion}
\subsection{Why REINFORCE performed well in this PdM setup}
The task has a small action space and ends on replacement, but the policy still has to infer timing from sensor trajectories. In this rerun, \textsc{REINFORCE} becomes the best mean performer once attention is enabled in both schemas (SIT mean with attention 0.7492, IEEE mean with attention 0.8569), whereas it does not beat \textsc{PPO} in baseline-only comparisons. This pattern indicates that representation quality is a primary driver for \textsc{REINFORCE} in this PdM setting.
The best mechanism is schema-dependent: \textsc{REINFORCE} peaks with Nadaraya--Watson on SIT but with self-attention on IEEE, while \textsc{PPO} peaks with multi-head on SIT and Nadaraya--Watson on IEEE (Table~\ref{tab:best_mech_ppo_rf}).

\subsection{Implications for industrial practitioners}
A practical deployment lesson is to benchmark no-attention first, then add attention selectively. For SIT-like schemas, attention can provide clear gains for strong policy-gradient learners. For IEEE-like schemas, pooled attention gains are not guaranteed even when some individual mechanisms improve selected learners. Therefore, encoder choice should be treated as schema-specific model selection, validated with significance testing rather than assumed to transfer.

\subsection{Relation to the literature gap}
The literature review suggests that most prior predictive-maintenance studies place attention in prognostics or state encoding rather than in a directly comparable RL study across multiple extractors. The present results fit that gap more closely than a pure RUL study: they test whether attention-style state representations materially change maintenance-policy performance under a common environment and evaluation workflow. At the same time, the results also preserve the caution from the literature. Attention is not universally beneficial, pooled gains are schema-dependent, and the strongest evidence here is for selected algorithm--schema pairings rather than for a blanket claim that attention dominates no-attention designs.

\section{Conclusion} \label{sec:conclusion} % 17-Jul-2026 retraining

\subsection{Summary of key findings}
This paper studied \textsc{REINFORCE} with attention-style feature extractors for milling-tool replacement. The literature review first clarified a narrower research gap than ``attention for PdM'' in general: the unresolved issue is how attention-based state representations affect RL maintenance performance, not merely upstream prognostics. Empirically, without attention, \textsc{PPO} is the best mean performer in both schemas. With attention, the best mean shifts to \textsc{REINFORCE} in both schemas. Focused Welch tests show that \textsc{REINFORCE} does not exceed \textsc{PPO} in baseline-only settings, but does so after adding attention. Mechanism-level synthesis shows schema-dependent winners: for \textsc{REINFORCE}, Nadaraya--Watson is best on SIT (0.8646 $\pm$ 0.2032) and self-attention is best on IEEE (0.8860 $\pm$ 0.0708); for \textsc{PPO}, multi-head is best on SIT (0.8004 $\pm$ 0.2694) and Nadaraya--Watson is best on IEEE (0.8343 $\pm$ 0.1867). The same environment design works across SIT and IEEE, and attention interoperability across schemas is strongest for \textsc{REINFORCE}, which improves significantly over its no-attention baseline in both SIT (0.4907 to 0.7492, $p=0.0000$) and IEEE (0.7496 to 0.8569, $p=0.0440$). This gain is not universal, however: pooled across all four algorithms, attention does not improve mean performance on IEEE (0.7377 to 0.5906, $p=1.0000$), so encoder selection must remain schema- and algorithm-specific.

\subsection{Future work}
Future work should test attention over short time windows, measure inference time and memory use, add richer maintenance costs, compare non-RL and lightweight no-attention encoders, and study robustness under sensor drift and cross-machine shift.

